\documentclass[public]{../../classes/nervtech-overview}

\nervsetup{
  company={NervTech},
  project={Node},
  document={Initial Project Overview},
  subtitle={A plain-language briefing for stakeholders, collaborators, and portfolio reviewers},
  docid={NT-N-OVR-000},
  revision={A},
  status={Draft Brief},
  author={Anhadh Virk},
  owner={NervTech Systems Engineering},
  date={2026-06-25},
  audience={Non-technical stakeholders; Portfolio reviewers},
  tagline={Node: The Building Block of NervTech Solutions},
  purpose={In order to produce more complex platforms to test reinforcement learning and advanced controls research, a robust, multipurpose actuator must be developed that has high precision, bandwith and torque capabilities. The Node fills this need, while being
           developed in house without the extra integration time and financial cost incurred by OEM integrated actuators.},
  summary={This overview will define the scope, goals and key deliverables, including timelines, of the Node project. This is essential in ensuring project success is met in a desirable timeframe.}
}

\begin{document}
\maketitle
\tableofcontents
\newpage

\section{Project Snapshot}

\projectsnapshot
{A compact, easily integratable robotic actuator module that combines motion generation, sensing, electronics, and control in one reusable unit.}
{Robotic systems become easier to build when the joint actuator is robust, reusable and self-contained in terms of control and sensing.}
{The project is in early architecture and prototype planning.}
{A 3D model of the first prototype in a simulated environment including the real life kinematics and dynamics of the system.}

\subsection*{Background}
As the labour automation industry becomes more expansive, older notions of developing
highly specialised robotic systems for specific industries are becoming obsolete.
As tasks with inherent uncertainty are being targeted for automation, distributed,
easily configurable robotic systems will form the bedrock of the next generation of
industry.
NervTech systems aims to be at the forefront of this movement, and the Node project
is a key step in that direction.
The project will develop a reusable, highly versatile and easily integratable robotic
actuator module that can be used in a variety of NervTech robotic systems.
The eventual goal will be to develop a family of robotic components, including multi-dof arms and end effectors, that can be easily configured to form a variety of robotic
systems for different applications.

\section{Problem Statement}
\textit{NervTech requires a compact, modular actuator that can provide highly precise controlled robotic motion for use in future distributed robotic platforms.
  The actuator must integrate sensing, control, communication, and mechanical mounting into a single reusable unit.
  Without such a platform, future robotic systems would require repeated custom integration of motors, sensors, drivers, and control hardware, increasing development time and reducing design consistency.}

\section{Goals}
\begin{goalstable}
  \goalrow{The actuator should provide high precision torque and position control in high motion frequency applications.}{The actuator should be able to move standard payloads up to 20kgs at frequencies comparable to human motion.}{Precision and Versatility}
  \goalrow{The actuator should be reusable and easily integratable into a variety of frames for robotic systems.}{The actuator should be able to be mounted and dismounted with conventional tools from any of the standard NervTech linkages and/or frames.}{Reusability and Versatility}
  \goalrow{The actuator should be able to provide a variety of sensing capabilities for advanced control.}{The actuator should be able to provide position, velocity, torque, and temperature sensing capabilities with minimal noise.}{Control and Versatility}
  \goalrow{The actuator should be able to provide a variety of communication capabilities with minimal latency.}{The actuator should be able to communicate with a variety of control systems using at least 2 standard industrial communication protocols, which can robustly deliver packets with sub-millisecond latency.}{Control and Communication}
  \goalrow{The actuator should be able to be powered by standard LV external DC sources.}{The actuator should be able to be powered via USB-C or standard molex connectors, from power supplies ranging from 24V - 48V.}{Versatility and Power}
  \goalrow{The actuator should include standard industrial safety features.}{The actuator should be able to provide safety features such as overcurrent protection, overtemperature protection, Safe Torque Off voltage lines and programmable max torque thresholds.}{Safety and Reliability}
  \goalrow{The actuator should be durable and reliable for long term use.}{The actuator should be able to operate for at least 10,000 hours of continuous operation without failure, and should be able to withstand standard industrial environmental conditions including dust and high moisture.}{Reliability and Durability}
\end{goalstable}

\section{Deliverables}
\begin{deliverablestable}
  \deliverablerow
  {NT-N-D1}{Initial project overview}
  {Plain-language document defining the product need, goals, requirements, success criteria, and roadmap.}
  {4/07/2026}

  \deliverablerow
  {NT-N-D2}{Preliminary research reports}
  {Comparison of actuator housing architectures, speed reducer architectures, power control architectures, communication protocols, safety feature implementations and motion controller designs.}
  {15/07/2026}

  \deliverablerow
  {NT-N-D3}{Technical design documents}
  {Mechanical, electrical, firmware and verification design documents, plus simulated environment using CAD, for the Node actuator.}
  {26/07/2026}

  \deliverablerow
  {NT-N-D4}{Prototype test report}
  {Bench-test evidence including motion response, repeatability, current draw, fault behaviour, and lessons learned.}
  {21/08/2026}
\end{deliverablestable}

\end{document}
